% Context from chapters/inverse-problems.tex; for mathematical reference, not compilation.
silon$. In this example, TV gives a modest visual improvement without increasing the SNR.

                                                                                                                                                                

\myfigure{
\tabdeux{
\image{invpbm}{.3}{inpainting-cameraman-image}&
\image{invpbm}{.3}{inpainting-cameraman-observations}\\
Image $f_0$ & Observation $y=\Phi f_0$ \\
\image{invpbm}{.3}{inpainting-cameraman-sobolev}&
\image{invpbm}{.3}{inpainting-cameraman-tv} \\
Sobolev $f^\star$ & TV $f^\star$
}
}{ 
	Inpainting with Sobolev and TV regularization.   
}{fig-inpainting-cameraman}



                                                                
\subsection{Tomography Inversion}

In an idealized two-dimensional tomography model, each projection integrates the unknown image along a line $\De_{t,\th}$ defined by
\eq{
	x \cdot \tau_\theta = x_1 \cos\th + x_2\sin\th = t
}
Here $\tau_\theta=(\cos\theta,\sin\theta)$ is the unit normal to the line.

The resulting Radon transform is
\eq{
	\foralls \th \in [0,\pi), \foralls t \in \RR, \quad
	p_{\th}(t) = \int_{\De_{t,\th}} f(x) \,d s
	 = \iint f(x)\, \de( x \cdot \tau_\theta - t )\, d x 
}
see Figure~\ref{fig-tomo-principle}.

\myfigure{
\image{invpbm}{.4}{tomo-principle}
}{ 
	Principle of tomography acquisition.   
}{fig-tomo-principle}

With $\hat f(\omega)=\int f(x)e^{-\imath x\cdot\omega}\,\mathrm{d}x$, the Fourier slice theorem identifies the one-dimensional Fourier transform of each projection with a radial slice of the two-dimensional Fourier transform
\eql{\label{eq-fourier-slice}
	\foralls \th \in [0,\pi)~,~ 
\foralls \xi \in \RR\quad
		\hat p_\th(\xi) = \hat f( \xi \cos \th, \xi \sin \th ).
}
Changing to polar coordinates in Fourier inversion gives the filtered-backprojection formula
\eq{
	f(x) = 
\frac{1}{2\pi} \int_{0}^\pi (p_\th \star h)(x \cdot \tau_\theta)\, d \th
}
with $\hat h(\xi) = |\xi|$.


\myfigure{
\tabtrois{
                                
\image{invpbm}{.3}{tomo-image}&
\image{invpbm}{.35}{tomo-radon-subsample}&
\image{invpbm}{.3}{tomo-radon-fourier}\\
Image $f$ & Radon sub-sampling & Fourier domain
}
}{ 
	Partial Fourier measurements.  
}{fig-tomo-radon-subsample}


We model acquisition at finitely many equally spaced orientations $\{\th_k=\pi k/K\}_{0\leq k<K}$ by the partial Radon transform
\eq{
	R f = ( p_{\th_k} )_{0 \leq k < K}.
}
By~\eqref{eq-fourier-slice}, knowing $R f$ is equivalent to knowing the Fourier transform of $f$ along the corresponding radial lines:
\eq{
	\{ \hat f(\xi \cos(\th_k), \xi \sin(\th_k)) \}_k.
}
For a simplified discrete model, we therefore treat acquisition as direct sampling of the Fourier transform:
\eq{
	\Phi f = ( \hat f[\om] )_{\om \in \Om} \in \CC^P
}
where $\Om$ consists of discrete radial sampling lines in the Fourier plane; see Figure \ref{fig-tomo-radon-subsample}, right.

In this idealized discrete model, recovery is an inpainting problem in the Fourier domain. We use the unitary discrete Fourier transform and consider measurements
\eq{
	\foralls \om \in \Om, \quad y[\om] = \hat f[\om] + w[\om]
}
where $w[\om]$ denotes measurement noise, modeled here as Gaussian white noise.

\myfigure{
\tabtrois{
\image{invpbm}{.3}{tomo-image}&
\image{invpbm}{.3}{tomo-pinv-13}&
\image{invpbm}{.3}{tomo-pinv-32}\\
Image $f_0$ & 13 projections & 32 projections.
}
}{ 
	Pseudoinverse reconstruction from partial Radon projections.  
}{fig-tomo-pinv}

For these partial Fourier measurements, the pseudoinverse $f^+ = \Phi^+ y$ from \eqref{eq-pseudo-inverse} is given by
\eq{
	{\hat f}^+[\om] = \choice{
		y[\om] \qifq \om \in \Om,\\
		\; 0 \qifq \om \notin \Om.
	}
}
Figure \ref{fig-tomo-pinv} shows examples of pseudoinverse reconstruction as the size of $\Om$ increases. This reconstruction exhibits pronounced artifacts because missing Fourier frequencies are set to zero.

The total variation regularization \eqref{eq-ip-regul} reads
\eq{
	f^\star \in \uargmin{f} \frac{1}{2}\sum_{\om \in \Om} |y[\om] - \hat f[\om]|^2 
	+ \la \normTV{f}.
}
TV regularization suits images with approximately constant regions separated by sharp boundaries, as in the cartoon model of Section \ref{subsec-cartoon-images}. Figure \ref{fig-tomo-tv} compares TV reconstruction and the pseudoinverse on a synthetic image of this type. The example illustrates recovery of sharp features from incomplete Fourier data. Spatial inpainting can be more challenging for TV, as Figure~\ref{fig-inpainting-cameraman} illustrates.

\myfigure{
\tabtrois{
\image{invpbm}{.3}{tomo-image}&
\image{invpbm}{.3}{tomo-pinv-fourier}&
\image{invpbm}{.3}{tomo-tv}\\
Image $f_0$ & Pseudoinverse & TV
}
}{ 
    Tomographic reconstruction with total variation regularization.  
}{fig-tomo-tv}
